% Generated draft: requires mathematical and source review.
\documentclass[11pt]{article}
\usepackage[margin=25mm]{geometry}
\usepackage{fontspec}
\setmainfont{Latin Modern Roman}
\usepackage{amsmath,amssymb,mathtools,unicode-math}
\setmathfont{Latin Modern Math}
\usepackage{enumitem,needspace,xurl,hyperref}
\setlength{\parindent}{0pt}
\setlength{\parskip}{5pt}
\setlength{\emergencystretch}{3em}
\newcommand{\N}{\mathbb{N}}
\newcommand{\Z}{\mathbb{Z}}
\newcommand{\Q}{\mathbb{Q}}
\newcommand{\R}{\mathbb{R}}
\newcommand{\C}{\mathbb{C}}
\newcommand{\F}{\mathbb{F}}
\newcommand{\A}{\mathbb{A}}
\newcommand{\PP}{\mathbb{P}}
\newcommand{\E}{\mathbb{E}}
\newcommand{\eps}{\varepsilon}
\newcommand{\dd}{\,\mathrm{d}}
\begin{document}
\section*{UnsolvedMath ID: 9400051}
\textbf{Source problem code:} catalogue entry
\subsection{Definitions and mathematical statement}
For a positive integer $n$, let
\[
A_n=\{x\in\R:x\text{ is algebraic over }\Q,\ [\Q(x):\Q]=n\}.
\]
A sequence $(a_j)_{j\geq0}$ is eventually periodic if there exist integers $N\geq0$ and $d\geq1$ such that $a_{j+d}=a_j$ for every $j\geq N$. The literal source asks, for every $n$, for an assignment
\[
F_n:\R\longrightarrow\N^{\N}
\]
whose output is eventually periodic exactly for $x\in A_n$. Degree is exactly $n$, not degree at most $n$, and the assignment may depend on $n$.

As written, the source imposes no algorithmic, dynamical, structural or naturality constraint on $F_n$. An arbitrary assignment may branch on membership in $A_n$, making this bare existence criterion fail to express the substantive Hermite problem. The supplied summary instead points to a natural continued-fraction-type expansion or multidimensional algorithm; its input/output rules and success criteria are not specified.

The degree-two benchmark is Lagrange's theorem: a nonterminating regular continued fraction is eventually periodic exactly for a real quadratic irrational. Rational numbers terminate under the usual convention. Extending every finite expansion by an eventually periodic tail would spoil an exact-degree-two criterion on all real numbers, so finite outputs or a rational-input encoding must be specified. The source's natural-number sequence formulation also needs any sign and initial-digit conventions. Thus that benchmark cannot silently be treated as a literal definition of $F_2$ on every real input.

The requested higher-degree result should provide the intended natural expansion with the exact periodicity characterization, rather than an unrestricted set-theoretic coding. No such algorithm or definition of natural is supplied.
\subsection{Short English statement}
Find a suitably natural expansion whose eventual periodicity characterizes real algebraic numbers of each exact degree. Without constraints on the assignment, the literal coding question is not the substantive Hermite problem.
\subsection{Sources}
\begin{itemize}
\item Local UnsolvedMath registry, permanent ID 9400051, exact source code catalogue entry, ``Hermite's problem''; no proposer or year supplied. Registry status partially solved, classification PARTIAL-PROGRESS.
\item Wikipedia contributors, \emph{List of unsolved problems in mathematics}, Algebraic number theory bullet 51, wikitext line 64, pinned revision 1366636767, \url{https://en.wikipedia.org/w/index.php?title=List_of_unsolved_problems_in_mathematics&oldid=1366636767#Number_theory}. Supplied MediaWiki extraction accessed July 29, 2026. The entry is staged pending semantic deduplication and primary-source review.
\item The supplied summary invokes Lagrange's continued-fraction theorem for degree two and describes the higher-degree natural multidimensional-expansion question as open, referring generally to Wikipedia's ``Hermite's problem'' without a primary citation or exact algorithmic formulation.
\item Provenance flags current primary-literature status and redistribution terms NEEDS\_REVIEW. Neither staging nor the registry published flag establishes source/status review.
\end{itemize}
\end{document}
